
\textbf{Proposition 3.2.} \uline{Soit} @@M0@@ \uline{un morphisme projectif}, \uline{d'intersection complète}, \uline{et dont la dimension relative virtuelle} @@M1@@ \uline{est constante}. \uline{On suppose que} @@M2@@ \uline{est quasi-compact}, \uline{et admet un faisceau ample}. \uline{On a alors pour tout} @@M3@@:
@@M4@@

On obtient donc en particulier un homomorphisme additif :
@@M5@@
envoyant @@M6@@ dans @@M7@@.

\textbf{3.3.} Rappelons que nous avons défini dans (V 6.1) un foncteur @@M8@@ sur la catégorie des @@M9@@-algèbres graduées, à valeurs dans la catégorie des @@M10@@-@@M11@@-algèbres augmentées, @@M12@@ étant un anneau binômial (V 2.7). Si @@M13@@ est une @@M14@@-algèbre graduée, on a :
@@M15@@

Pour toute @@M16@@-@@M17@@-algèbre augmentée @@M18@@, nous avons défini un @@M19@@-@@M20@@-homomorphisme, nommé «~classe de Chern complétée~» (V 6.7):
@@M21@@

Enfin, nous avons défini un homomorphisme d'anneaux, fonctoriel en @@M22@@, et nommé caractère de Chern (V 6.4):
@@M23@@
et un endomorphisme de groupes, nommé caractère de Todd (V 1.16):
@@M24@@
